\section{Introduction}
\IEEEPARstart{P}{artitioned} time stepping is a practical necessity when simulating stiff, multi-physics robotic systems under real-time constraints. By distributing subsystem dynamics across parallel computational units, model partitioning can substantially reduce wall-clock time per macro-step \cite{negrut2014parallel, schweizer2015explicit}.

Once models are split, the bottleneck becomes the numerical interface. Discrete exchange can break power consistency and create coupling defects \cite{sadjina2024energy,sadjina2017energy,benedikt2013nepce}, yielding spurious energy injection (effective negative damping) even when subsystems are passive \cite{sadjina2024energy,gonzalez2019energy,rodriguez2022energy}. 
An interface must therefore enforce discrete passivity of the coupled macro-step.

Related work on partitioned interfaces spans non-iterative (explicit) and iterative schemes \cite{gomes2018co,blochwitz2011functional,schweizer2015explicit}: non-iterative zero-order-hold or extrapolation is attractive for real time but can drift or destabilize under stiff coupling \cite{peiret2018multibody,dai2023model,raoofian2024nonsmooth}, while macro-step iterations (Gauss--Seidel/Jacobi, relaxation) improve consistency \cite{dai2023model,raoofian2024nonsmooth} but raise delicate finite-iteration stability questions \cite{schweizer2016implicit}. 
Transmission Line Modeling (TLM) introduces physically motivated time delays to break algebraic loops and preserve passivity \cite{krus2011robust,braun2013multi,braun2022numerically}. Fixed delays can distort high-frequency behavior and limit performance \cite{barbierato2024locally,burton1993analysis,braun2024transmission}, motivating delay-free alternatives or impedance tuning.

Structure-preserving modeling makes the interface an explicit power contract. Port-Hamiltonian formulations encode interconnection as a power-conserving constraint \cite{celledoni2017energy,parks2017variational,ehrhardt2025structure}, and structure-preserving time stepping is well developed for monolithic integration \cite{celledoni2017energy,argus2021theory,ehrhardt2025structure} but harder to retain under partitioned, iterative coupling, notably for DAEs and pHDAEs \cite{bartel2023operator,ehrhardt2025structure}. In port coordinates the coupling layer is an input–output map whose inequality quantifies equivalent damping. This motivates energy-safe interface updates via operator-splitting compositions \cite{bartel2023operator,bartel2025splitting}, whose guarantees are typically asymptotic, so finite intermediate iterates may still fail to remain physically admissible or passive.

Wave (scattering) variables turn power coupling into passive message passing \cite{cervera2007interconnection} and are closely related to TLM interfaces when explicit delays are introduced \cite{krus2011robust,braun2013multi,braun2022numerically,braun2024transmission}. 
Yet stiff delay-free algebraic loops often still require interface tuning in practice \cite{chen2021explicit,schweizer2015explicit}, while operator-splitting methods provide asymptotic consistency guarantees but do not by themselves yield a wave-domain, iterate-level energy accounting. 
This motivates embedding splitting updates in the wave domain to obtain parallel coupling with an energy guarantee under finite inner iterations.

In this letter, we propose an \emph{early-terminable} energy-safe coupling interface by embedding operator-splitting updates in wave (scattering) coordinates. 
At each macro-step, a Douglas--Rachford (DR) inner iteration reconciles a linear lossless wave coupling constraint with the discrete-time port behavior induced by the subsystem integrators. The inner-iteration budget $K_n$ provides a tunable accuracy--effort trade-off, together with an augmented-storage energy guarantee under any finite $K_n$.
The main contributions are as follows:

1) \textbf{An early-terminable energy-safe interface for parallel coupling:}
We introduce a parallelizable iterative coupling interface in wave (scattering) variables that enforces a linear lossless interconnection via Douglas--Rachford (DR) inner iterations. The method applies to general nonlinear subsystems and remains energy-safe for any finite inner-iteration budget (Algorithm~\ref{alg:scattering_iterative_coupling}).

2) \textbf{Finite-iteration passivity guarantee:}
For any finite inner-iteration budget, we establish a discrete passivity guarantee for the macro-step update via an augmented storage function (Theorem~\ref{thm:augmented_storage}).

3) \textbf{Convergence to the monolithic discretization:}
As the inner-iteration budget increases, the partitioned update provably converges to the monolithic discrete-time update induced by the same integrator (Theorem~\ref{thm:kn_to_infty_monolithic}).


The remainder of the letter is organized as follows: Section~II models the subsystems and coupling in the wave domain and formulates the interface as a fixed-point problem; Section~III presents the iterative coupling algorithm; Section~IV establishes passivity and convergence properties; Section~V reports numerical experiments; and Section~VI concludes.


